(1) 无流无耗散声压满足 \(p_{tt}=c^2\nabla^2p\)，刚性壁要求流体法向速度 \(u_r(R,t)=0\)。由线性Euler方程 \(\rho u_{r,t}=-\partial_rp\)，对非零频率谐振声场可等价写 \(\partial_rp(R,t)=0\)，即法向Neumann边界。
(2) 令 \(p(r,t)=P(r)e^{-i\omega t}\)，\(k=\omega/c\)，径向方程为 \(P''+2P'/r+k^2P=0\)。一般解为 \(A\sin(kr)/(kr)+B\cos(kr)/(kr)\)；后者在 \(r\to0\) 发散，球心正则性要求 \(B=0\)，故 \(P(r)=A j_0(kr)\)，\(j_0(x)=\sin x/x\)。
(3) 在刚壁 \(\partial_rP(R)=0\)，非零波数满足 \(j_0'(kR)=0\)，即 \(kR\cos(kR)-\sin(kR)=0\)，所以 \(\tan(kR)=kR\)（需排除 \(kR=0\) 的平凡除法情形）。第一个严格正根为 \(x_1\simeq4.493409\)，于是 \(\omega_1=cx_1/R\)。
(4) 球腔的非零径向频率按 \(\omega_n=c x_n/R\) 缩放。\(k=0\) 对应空间均匀压强的静态或极慢整体状态改变，非零频率下它不构成独立传播声振动模；刚壁限制的是法向粒子速度而非压强本身，若误设 \(P(R)=0\) 就会算出另一套错误的Dirichlet根。
